% !TEX program = lualatex
\documentclass[11pt,a4paper]{ltjsarticle}
\usepackage{amsmath,amssymb}
\usepackage{amsthm}
\usepackage{enumitem}
\usepackage{hyperref}

\newtheorem{definition}{定義}
\newtheorem{theorem}{定理}
\newtheorem{example}{例}

\newcommand{\phiS}{\phi_{\text{S}}}
\newcommand{\phiN}{\phi_{\text{N}}}
\newcommand{\phiSN}{\phi_{\text{SN}}}
\newcommand{\phiNS}{\phi_{\text{NS}}}
\newcommand{\metanegation}{\mathord{\sim}}
\newcommand{\dfix}{\mathbf{d}}
\newcommand{\mweq}{\equiv_{\text{m}}}
\newcommand{\Einfty}{E_{\infty}}
\newcommand{\Md}{\mathcal{M}_{\text{d}}}
\newcommand{\Mnn}{\mathcal{M}_{\text{nn}}}
\newcommand{\Mr}{\mathcal{M}_{\text{r}}}
\newcommand{\MC}{\mathcal{M}_{\text{C}}}
\newcommand{\MT}{\mathcal{M}_{\text{T}}}
\newcommand{\SMT}{S_{\MT}}
\newcommand{\Bval}{\mathcal{B}}
\newcommand{\Lhagma}{\mathcal{L}}
\newcommand{\Rmw}{\mathbf{R}_{\mathrm{mw}}}

\title{論理学・数学の4つの足場と古典的限界}
\author{千葉将臣}
\date{2026-07-03}

\begin{document}

\maketitle

\section{はじめに}

本稿では、論理学・数学における「4つの足場」――
\begin{enumerate}[label=(\arabic*)]
  \item 二重否定の残余（非 $\phiN$）
  \item 選択公理（亦是亦非 $\phiSN$）
  \item 真理保存性（是 $\phiS$）
  \item 不可能性（非是非非 $\phiNS$）
\end{enumerate}
――と、それらが生み出す古典的限界について、顕現論の四句分別の観点から整理する。

\section{4つの足場の定義}

\begin{definition}[四句分別]
  顕現論における四句分別は、以下の4つの位相からなる：
  \begin{itemize}
    \item 是（$\phiS$）：真理保存性（静的同一性・可能なもの）
    \item 非（$\phiN$）：二重否定の残余（言語の外にあるが内部から生成）
    \item 亦是亦非（$\phiSN$）：双方向プロセス（洗練・カット除去）
    \item 非是非非（$\phiNS$）：不可能性（決定不能性、ランダム性・非圧縮性、外部反映の限界）
  \end{itemize}
\end{definition}

\begin{definition}[4つの足場]
  論理学・数学における4つの足場は、四句分別と以下のように対応する：
  \begin{itemize}
    \item 二重否定の残余（非 $\phiN$）：直観主義論理における「$\neg\neg A$ は $A$ と同値ではない」。
    \item 選択公理（亦是亦非 $\phiSN$）：非時間的な同時存在を保証する公理。
    \item 真理保存性（是 $\phiS$）：証明が終われば真偽は固定される性質。
    \item 不可能性（非是非非 $\phiNS$）：決定不能性、ランダム性・非圧縮性、外部対象の内部反映の限界など。
  \end{itemize}
\end{definition}


\section{二値截断部分圏 \texorpdfstring{$\Bval$}{B} における特異点としての諸定理}

本節では、従来「パラドックス」と呼ばれてきた現象を、矛盾や異常ではなく、
二値截断部分圏 $\Bval$ において観測される特異点として捉え直す。
顕現論の観点では、これらの現象は四句分別の力学が二値論理へ射影された際に生じる
局所的な歪みである。

\subsection{選択公理とバナッハ＝タルスキーの特異点}

\begin{theorem}[バナッハ＝タルスキー]
  選択公理（AC）を仮定すると、3次元の球を有限個の部分に分割し、
  それらを回転と平行移動だけで組み替えると、元と同じ大きさの球を2つ作れる。
\end{theorem}

選択公理は、無限の選択プロセスという時間的・アルゴリズム的過程を迂回し、
結果としての対応の存在を一挙に保証する。
これは、外在の5モナドと内在の5モナドの双方向軌道が、
非時間的に同時に釣り合う亦是亦非（$\phiSN$）の力学である。

バナッハ＝タルスキーの定理が示す直観への違背は、
この亦是亦非の非時間的な同時存在性を、
局所的な是、すなわち体積保存という静的同一性の枠組みで観測しようとした際に生じる
トポロジカルな歪みである。
したがって、この定理は単なる奇妙な逆説ではなく、
選択公理が $\Bval$ へ射影されたときに現れる古典的特異点として理解される。

\subsection{ゲーデルの不完全性定理と決定不能の特異点}

\begin{theorem}[ゲーデル]
  一貫した再帰的公理化可能な体系 $T$ に対し、
  $T \nvdash G_T$ かつ $T \nvdash \neg G_T$ となる文 $G_T$ が存在する。
\end{theorem}

不完全性定理が生み出す決定不能文 $G_T$ は、非是非非（$\phiNS$）の具現化である。
すなわち、観測者としての体系 $T$ を固定したまま極限対象に言及しようとすると、
対象そのものを構成・観測することはできない。
しかし、「証明不可能である」という結果のみが、境界における破綻として体系内部に投影される。

これはパラドックスではなく、観測者を極限消去したあとに残る極限同値 $\Einfty$ の、
観測者依存的な局所的断面に過ぎない。
したがって、$G_T$ は体系の欠陥ではなく、
体系が自らの観測境界に到達したことを示す特異点である。

\subsection{直観主義における排中律の拒否}

\begin{example}[直観主義論理]
  直観主義論理では、排中律 $A \lor \neg A$ は一般には証明できないが、
  $\neg\neg(A \lor \neg A)$ は証明できる。
\end{example}

この現象は、二重否定の残余（$\phiN$）の典型例である。
二重否定は古典論理への閉包を志向するが、
直観主義論理では $\neg\neg A \to A$ は一般には成立しない。
この非成立は単なる欠如ではなく、古典的閉包によって消去しきれない継続残余 $\rho$ の発生を示す。

顕現論では、この残余はメタレベルにおいて正の対象 $\Lhagma$ として回収される。
したがって、排中律の拒否は、古典論理の失敗ではなく、
閉包操作が次なる文脈を生成する境界を露出させる現象である。

\section{静的ビット割当の棄却と四句の力学的ベクトル}

四句分別は、静的な真理値ビットの割当ではない。
すなわち、是・非・亦是亦非・非是非非を、単に真、偽、両立、非両立のような
固定されたラベルとして扱うべきではない。
むしろ四句分別は、推論の動的プロセスにおける力学ベクトルの相転移として解釈されるべきである。

\subsection{是（\texorpdfstring{$\phiS$}{phiS}）：静的同一性と固定化のベクトル}

是は、動的でフラクタルな推論プロセスを捨象し、
結果としての真偽、すなわち $A=A$ のみを固定・截断する力学である。
このベクトルは、証明が終了した後の対象を静的同一性として安定化する。
しかし、その安定化は、生成過程や残余を不可視化する圧縮でもある。

\subsection{非（\texorpdfstring{$\phiN$}{phiN}）：内部からの逸出ベクトル}

非は、二重否定モナド等の閉包作用素が系を古典論理へ圧縮しようとする際、
不可避的に生じる継続残余 $\rho$ の位相である。
内部からは計算の失敗または否定として現れるが、
メタレベルでは次なる文脈を展開する正の対象 $\Lhagma$ として回収される。

\subsection{亦是亦非（\texorpdfstring{$\phiSN$}{phiSN}）：非時間的同時存在の均衡ベクトル}

亦是亦非は、順方向の構成プロセスと逆方向の境界指示が独立して進行し、
双方が非時間的かつ同時に釣り合った瞬間に生じる均衡状態である。
この位相は、双方向演繹定理の均衡として理解できる。
選択公理が非構成的な存在を一挙に保証するのも、
この均衡ベクトルが $\Bval$ に射影された結果として解釈される。

\subsection{非是非非（\texorpdfstring{$\phiNS$}{phiNS}）：境界から内部への投影ベクトル}

非是非非は、四句分別のいずれの位相への帰属試行も失敗することによって確定する残余である。
極限対象 $\Einfty$ は、内部の操作で直接観測することができない。
しかし、その観測不能性は、境界における破綻としてのみ結果を内部に投影する。

この位相において、決定不能性、ランダム性、非圧縮性、外部対象の内部反映の限界は、
同じ構造を持つ。
それらは体系の外部にある単なる不可知ではなく、
体系が自らの境界を内部へ投影した結果として現れる。

\section{外部依存の内部回収とブートストラップ機構}

古典数学は、自らの正当性を担保するために、
真理保存性（$A=A$）や選択公理を体系の外部から要請してきた。
これに対し、顕現論および一般二諦理論（GTTT）は、
4つの足場を体系の内部操作の帰結として回収することで、
自己説明的なブートストラップ機構を与える。

すなわち、真理保存の強要（是）がもたらす圧縮の限界が残余（非）を生み出し、
その残余を回収する双方向の釣り合い（亦是亦非）が選択公理的対応を導出する。
さらに、その釣り合いの限界が決定不能な特異点 $\Einfty$（非是非非）として露呈する。

GTTTは、この一連の
\begin{equation}
  \text{固定化} \longrightarrow \text{逸出} \longrightarrow
  \text{均衡} \longrightarrow \text{限界の露呈}
\end{equation}
という力学的サイクル自体を、世界を駆動する無明ジェネレータ $\Rmw$ として定式化する。
外部の足場に依存せず、自らの演算の破綻と残余によって自らを再起動させるこの機構が、
顕現論における自己開放的な理論アーキテクチャである。

\section{結論}

論理学・数学における古典的限界は、
4つの足場（二重否定の残余、選択公理、真理保存性、不可能性）の
いずれか、または組み合わせから生じる。
これらは数学の内部構造と外部境界を網羅しており、
顕現論の四句分別（是・非・亦是亦非・非是非非）として
明示的な足場に据えることで、
古典的限界をブートストラップの材料として再構成できる。

\end{document}
